% =====================================================================
%  Additional explanatory figures for v2 (vector only, TikZ)
%  Styles fbox / arr / alab are defined in figures-old.tex.
% =====================================================================

% ---------------------------------------------------------------------
% X1 — the three halves feeding Clay
% ---------------------------------------------------------------------
\newcommand{\figthreehalves}{%
\begin{figure}[t]
\centering
\begin{tikzpicture}[x=1cm,y=1cm]
  \node[fbox, fill=cC!8, draw=cC!55, text width=4.3cm, minimum height=2.7cm] (g1) at (2.7,3.5)
    {\textbf{G1 --- existence}\\[2pt]
     \scriptsize 4D continuum, OS0--OS4;\\
     fixed volume assembled\\
     modulo the uniform $Z_a$\\[2pt]
     \LT{} mod.\ L3 \quad \LM{} L3 (UV)};
  \node[fbox, fill=cM!8, draw=cM!60, text width=4.3cm, minimum height=2.7cm] (g2) at (8.0,3.5)
    {\textbf{G2 --- the gap}\\[2pt]
     \scriptsize $f_1^2>0$ with isolation;\\
     equivalently $G''(0)\ge\lambda$\\
     as $k\to0$, $(L,a)$-uniform\\[2pt]
     \LT{} region $\sim$90\% \quad \LM{} sign 0\%};
  \node[fbox, fill=cT!8, draw=cT!55, text width=4.3cm, minimum height=2.7cm] (ch) at (13.3,3.5)
    {\textbf{The channel}\\[2pt]
     \scriptsize lightest state in $0^{++}$;\\
     no-dissolution theorem;\\
     ordering input \code{IN-11}\\[2pt]
     \LC{} conditional};
  \node[fbox, fill=cN!7, draw=cN!55, text width=7.4cm] (clay) at (8.0,0.55)
    {\textbf{Clay: existence $\wedge$ uniform gap}
     \quad\textbf{0\% --- open}};

  \draw[arr] (g1.south) -- (clay.north west);
  \draw[arr] (g2.south) -- (clay.north);
  \draw[arr] (ch.south) -- (clay.north east);
\end{tikzpicture}
\caption{The three objects the campaign separates. G1 is a construction problem, G2 is a lower
bound on the Hamiltonian, and the channel is the identification of the lightest state with the
$0^{++}$ glueball. They feed one conclusion but are neither equivalent nor ordered: a gap does
not give existence, existence does not give a gap, and a gap in some channel does not identify
the channel. Only the middle object is undecided on both its region and its sign.}
\label{fig:threehalves}
\end{figure}
}

% ---------------------------------------------------------------------
% X2 — the research loop
% ---------------------------------------------------------------------
\newcommand{\figloop}{%
\begin{figure}[t]
\centering
\begin{tikzpicture}[x=1cm,y=1cm]
  \node[fbox, fill=cC!8, draw=cC!55, text width=3.6cm] (forge) at (8,5.6)
    {\textbf{1. Forge}\\ a new tool or identity};
  \node[fbox, fill=cE!8, draw=cE!55, text width=3.6cm] (acid) at (12.0,2.9)
    {\textbf{2. Acid test}\\ must fail on the free field};
  \node[fbox, fill=cT!8, draw=cT!55, text width=3.6cm] (audit) at (8,0.2)
    {\textbf{3. Audit}\\ against the source files};
  \node[fbox, fill=cM!8, draw=cM!60, text width=3.6cm] (bank) at (4.0,2.9)
    {\textbf{4. Bank or record}\\ theorem, or exact no-go};

  \draw[arr] (forge) -- node[alab, above right=-1pt] {\scriptsize test} (acid);
  \draw[arr] (acid) -- node[alab, below right=-1pt] {\scriptsize if it survives} (audit);
  \draw[arr] (audit) -- node[alab, below left=-1pt] {\scriptsize verdict} (bank);
  \draw[arr] (bank) -- node[alab, above left=-1pt] {\scriptsize next round} (forge);
  \node[font=\scriptsize, text=black!60, align=center, text width=5.4cm] at (8,2.9)
    {one round\\[2pt] a tool that passes\\ the free field is vacuous};
\end{tikzpicture}
\caption{The working loop imposed by the project charter. Every proposed tool must first fail on
the massless free field (where $f_1^2=0$ and $G''=0$); a tool that ``proves'' a gap there is
vacuous and is discarded. Only then is it run against the campaign's source files, and only then
is it banked as a theorem or recorded as a no-go with the exact point of failure.}
\label{fig:loop}
\end{figure}
}

% ---------------------------------------------------------------------
% X3 — the five labels
% ---------------------------------------------------------------------
\newcommand{\figlabelmap}{%
\begin{figure}[t]
\centering
\begin{tikzpicture}[x=1cm,y=1cm,font=\scriptsize]
  \foreach \x/\col/\chip/\name/\txt in {%
    1.8/cT/\LT/{Theorem}/{proved; imported results are named},
    5.0/cC/\LC/{Conditional}/{proved implication from named hypotheses},
    8.2/cN/\LN/{No-go}/{a whole class cannot work},
    11.4/cM/\LM/{Missing step}/{one open inequality, no slack hidden},
    14.6/cE/\LE/{Evidence}/{measured, not proved}}{%
    \node[fbox, fill=\col!6, draw=\col!45, text width=2.9cm, minimum height=2.5cm]
      at (\x,1.6) {\chip\\[2pt]\textbf{\name}\\[2pt]\txt};}
\end{tikzpicture}
\caption{The five labels used throughout. They are not decoration: each claim carries exactly
one, and the label decides what the claim may be used for. In particular an \LE{} number is never
substituted into a \LT{} statement, and a \LM{} is stated as an inequality with zero slack rather
than as ``work in progress''.}
\label{fig:labelmap}
\end{figure}
}

% ---------------------------------------------------------------------
% X4 — the spectrum ladder
% ---------------------------------------------------------------------
\newcommand{\figladder}{%
\begin{figure}[t]
\centering
\begin{tikzpicture}[x=1cm,y=1cm,font=\scriptsize]
  \draw[arr] (1.0,1.2) -- (15.2,1.2) node[alab, right] {$E$};
  \fill[gray!18] (12.0,1.2) rectangle (15.0,4.2);
  \draw[dashed, black!55] (12.0,1.2) -- (12.0,4.2);
  % levels
  \draw[very thick, black!70] (2.6,1.2) -- (2.6,3.0);
  \node[above] at (2.6,3.0) {$\Omega$ (vacuum)};
  \draw[very thick, cN!85] (6.0,1.2) -- (6.0,3.8);
  \node[above, text=cN!85] at (6.0,3.8) {$0^{++}$ at $m_1$};
  \node[left, text=cN!85] at (5.9,2.6) {height $f_1^2=\nu(\{m_1^2\})$};
  \draw[very thick, black!45] (9.0,1.2) -- (9.0,3.0);
  \node[above, text=black!55] at (9.0,3.0) {$0^{++*}$ at $\approx1.5\,m_1$};
  \node[text=black!60] at (13.5,3.95) {two-glueball};
  \node[text=black!60] at (13.5,3.65) {continuum};
  % brackets
  \draw[decorate, decoration={brace, amplitude=4pt}, cM!85]
    (6.15,2.2) -- node[above=4pt, text=cM!85] {$I_2$ margin $\delta\approx0.5\,m_1$} (8.85,2.2);
  \draw[decorate, decoration={brace, amplitude=4pt, mirror}, black!55]
    (6.15,0.95) -- node[below=4pt] {two-body threshold $2m_1$} (11.85,0.95);
\end{tikzpicture}
\caption{The spectral ladder of the $0^{++}$ channel. The atom weight $f_1^2$ (label I1) is the
height of the first level; the isolation (I2) is its distance to the \emph{excited} scalar, not
to the two-glueball continuum. This distinction is the reason the strong-coupling residue theorem
does not by itself give a gap, and why the operative threshold in the no-dissolution theorem is
$1.5\,m_g$ rather than $2m_g$.}
\label{fig:ladder}
\end{figure}
}

% ---------------------------------------------------------------------
% X5 — the G1 pipeline
% ---------------------------------------------------------------------
\newcommand{\figpipeline}{%
\begin{figure}[t]
\centering
\begin{tikzpicture}[x=1cm,y=1cm,font=\scriptsize]
  \foreach \x/\col/\nm/\txt in {%
    1.6/cT/{L1}/{definition of the measure},
    4.2/cT/{L2}/{single-block bound},
    6.8/cM/{L3}/{uniform $Z_a$},
    9.4/cT/{GL3}/{compactness},
    12.0/cT/{GL4}/{reflection positivity},
    14.5/cT/{OS0--4}/{reconstruction}}{%
    \node[fbox, fill=\col!7, draw=\col!50, text width=2.25cm, minimum height=1.7cm]
      at (\x,1.8) {\textbf{\nm}\\[1pt]\txt};}
  \foreach \a/\b in {1.6/4.2,4.2/6.8,6.8/9.4,9.4/12.0,12.0/14.5}{%
    \draw[arr] (\a+1.15,1.8) -- (\b-1.15,1.8);}
  \node[font=\scriptsize, text=cM!85, align=center, text width=5.4cm] at (6.8,-0.15)
    {the only open UV node: $\sup_a Z_a^{-1}\lVert H_a\rVert_D<\infty$\\ (a single-sided force-variance bound)};
\end{tikzpicture}
\caption{The existence pipeline G1. Every arrow except the third is closed by a theorem or by the
statement of the measure; L3 is the single open ultraviolet node. It is deliberately drawn next
to G2's H1 in the text, because L3 and H1 are the same species of inequality at two different
scales.}
\label{fig:pipeline}
\end{figure}
}

% ---------------------------------------------------------------------
% X6 — isometric lattice with a plaquette and a flux tube
% ---------------------------------------------------------------------
\newcommand{\figlattice}{%
\begin{figure}[t]
\centering
\begin{tikzpicture}[x=1cm,y=1cm,font=\scriptsize,
  X/.style={draw=black!45, line width=0.5pt},
  Xhot/.style={draw=cM!90, line width=2.2pt}]
  % isometric basis
  \def\ex{0.62} \def\ey{0.36}
  \def\fz{0.72}
  % vertices (i,j,k) in 0..3
  \foreach \i in {0,1,2,3}{%
    \foreach \j in {0,1,2,3}{%
      \foreach \k in {0,1,2,3}{%
        \pgfmathsetmacro{\px}{\i*\ex - \j*\ex}
        \pgfmathsetmacro{\py}{\i*\ey + \j*\ey + \k*\fz}
        \fill[black!70] (\px,\py) circle (1.1pt);
      }}}
  % links along x, y, z (only within 0..2)
  \foreach \i in {0,1,2}{%
    \foreach \j in {0,1,2}{%
      \foreach \k in {0,1,2}{%
        \pgfmathsetmacro{\ax}{\i*\ex-\j*\ex} \pgfmathsetmacro{\ay}{\i*\ey+\j*\ey+\k*\fz}
        \pgfmathsetmacro{\bx}{\ax+\ex} \pgfmathsetmacro{\by}{\ay+\ey}
        \draw[X] (\ax,\ay) -- (\bx,\by);
        \pgfmathsetmacro{\cx}{\ax-\ex} \pgfmathsetmacro{\cy}{\ay+\ey}
        \draw[X] (\ax,\ay) -- (\cx,\cy);
        \pgfmathsetmacro{\dx}{\ax} \pgfmathsetmacro{\dy}{\ay+\fz}
        \draw[X] (\ax,\ay) -- (\dx,\dy);
      }}}
  % highlighted plaquette on the top layer (i=0..1, j=0..1, k=3)
  \pgfmathsetmacro{\pax}{0} \pgfmathsetmacro{\pay}{0+0+3*\fz}
  \pgfmathsetmacro{\pbx}{\ex} \pgfmathsetmacro{\pby}{\ey+3*\fz}
  \pgfmathsetmacro{\pcx}{0} \pgfmathsetmacro{\pcy}{2*\ey+3*\fz}
  \pgfmathsetmacro{\pdx}{-\ex} \pgfmathsetmacro{\pdy}{\ey+3*\fz}
  \draw[Xhot] (\pax,\pay) -- (\pbx,\pby) -- (\pcx,\pcy) -- (\pdx,\pdy) -- cycle;
  % flux tube through the centre
  \draw[cN!80, line width=4pt, opacity=0.75] (-0.15,0.9) -- (-0.15,3.55);
  \node[text=cN!80, anchor=east] at (-0.45,2.1) {flux tube};
  \draw[cN!60, ->, thin] (-0.4,2.15) -- (-0.18,2.25);
  \node[text=cM!90] at (1.55,3.55) {plaquette $U_p$};
  \draw[cM!70, ->, thin] (1.15,3.45) -- (0.55,3.0);
\end{tikzpicture}
\caption{A $3^3$ periodic lattice in isometric view: site variables, link variables, one
plaquette $U_p$ (the object whose trace enters the Wilson action and the operator
$\operatorname{tr}F^2$), and a flux tube. The figure fixes the objects used throughout: the
strong-coupling residue theorem counts polymers of such plaquettes, and the no-dissolution
argument is about what a tube of this kind can and cannot do.}
\label{fig:lattice}
\end{figure}
}

% ---------------------------------------------------------------------
% X7 — the variance ledger
% ---------------------------------------------------------------------
\newcommand{\figbalance}{%
\begin{figure}[t]
\centering
\begin{tikzpicture}[x=1cm,y=1cm,font=\scriptsize]
  \draw[arr] (0.8,0.4) -- (0.8,4.6) node[alab, above] {value};
  \draw[dashed, black!45] (0.8,1.1) -- (14.6,1.1);
  % waterfall: E[H] up, Var down, W up, net
  \fill[cT!55] (1.7,1.1) rectangle (3.3,3.3);
  \node[above] at (2.5,3.3) {$E[H]$};
  \fill[cN!45] (4.4,1.3) rectangle (6.0,3.3);
  \node[below] at (5.2,1.3) {$-\operatorname{Var}$};
  \fill[cC!55] (7.1,1.1) rectangle (8.7,1.45);
  \node[above] at (7.9,1.45) {$+W_{\mathrm{full}}$};
  \fill[black!70] (9.8,1.1) rectangle (11.4,1.55);
  \node[above] at (10.6,1.55) {$G''=S-D$};
  \draw[dotted, black!60] (3.3,3.3) -- (4.4,3.3);
  \draw[dotted, black!60] (6.0,1.3) -- (7.1,1.3);
  \draw[dotted, black!60] (8.7,1.45) -- (9.8,1.45);
  \node[align=left, text width=3.4cm] at (13.2,2.6)
    {$G''=E[H]-\operatorname{Var}+W_{\mathrm{full}}$\\[2pt]
     $=\Delta_{\mathrm{dyn}}+W_{\mathrm{geo}}+X$};
  \node[text=black!55] at (5.2,0.55) {schematic heights};
\end{tikzpicture}
\caption{The variance ledger behind the gap. The unconditioned curvature $E[H]$ is large and
positive; the variance of the force subtracts from it; the geometric/cross remainder
$W_{\mathrm{full}}=W_{\mathrm{geo}}+X$ is what is left of the conditioning correction. The sign
of $G''$ is therefore not a statement about the size of $E[H]$ --- that is positive for trivial
reasons --- but about the subtractive variance term, which is exactly the object no known tool
bounds from above at $\beta\approx2.5$. Heights are schematic.}
\label{fig:balance}
\end{figure}
}

% ---------------------------------------------------------------------
% X8 — stable and unstable directions at the Savvidy point
% ---------------------------------------------------------------------
\newcommand{\figsaddle}{%
\begin{figure}[t]
\centering
\begin{tikzpicture}[x=1cm,y=1cm,font=\scriptsize]
  \draw[arr] (0.9,2.0) -- (7.2,2.0) node[alab, right] {$t$};
  \draw[arr] (0.9,0.4) -- (0.9,3.6) node[alab, above] {$V$};
  \draw[cN!85, thick, domain=-2.2:2.2, samples=60] plot ({3.6+1.25*\x},{2.0-0.28*\x*\x});
  \node[text=cN!85, align=left, text width=3.0cm] at (6.0,3.0)
    {negative mode:\\ $\lambda^2<0$ (unstable)};
  \draw[cT!70, thick, domain=-2.2:2.2, samples=60] plot ({11.0+1.25*\x},{2.0+0.28*\x*\x});
  \node[text=cT!70, align=left, text width=3.0cm] at (12.5,0.7)
    {positive modes:\\ stable};
  \node[text=black!60, align=center, text width=4.2cm] at (3.0,0.55)
    {transverse magnetic mode};
  \node[text=black!60, align=center, text width=4.2cm] at (12.6,3.5)
    {would-be convex directions};
\end{tikzpicture}
\caption{The Savvidy phenomenon in one dimension. Around a constant chromomagnetic background the
effective potential is concave along the transverse magnetic mode and convex along the others.
The concavity is physical, not a gauge artifact: it survives the gauge quotient, and it is the
reason every convexity-based tool (Pr\'ekopa, Brascamp--Lieb, log-concavity) stops at the
background. The campaign's $W_{\mathrm{geo}}\ge0$ theorem controls the Haar part, not this term.}
\label{fig:saddle}
\end{figure}
}

% ---------------------------------------------------------------------
% X9 — the theorem dependency graph
% ---------------------------------------------------------------------
\newcommand{\figproofdag}{%
\begin{figure}[t]
\centering
\begin{tikzpicture}[x=1cm,y=1cm,font=\scriptsize]
  \foreach \x/\col/\nm/\lab in {%
    1.7/cM/{H1}/\LM,
    4.9/cC/{CP}/\LC,
    8.1/cT/{$W_{\mathrm{geo}}\!\ge0$}/\LT,
    11.3/cT/{P3}/\LT,
    14.5/cC/{no-diss.}/\LC}{%
    \node[fbox, fill=\col!7, draw=\col!50, text width=2.7cm] at (\x,4.2)
      {\textbf{\nm}\\[1pt]\lab};}
  \node[fbox, fill=cT!10, draw=cT!60, text width=9.6cm] (red) at (8.0,1.9)
    {\textbf{conditional reduction (Lean 4)}\\[2pt]
     \code{clay\_gap\_conditional} : H1 $\wedge$ CP $\wedge$ $W_{\mathrm{geo}}\!\ge0$ $\wedge$ P3
     $\wedge$ no-diss. $\Longrightarrow$ $\Delta>0$\\[2pt]
     zero \code{sorryAx}; the five inputs are explicit hypotheses};
  \node[fbox, fill=cN!8, draw=cN!55, text width=9.6cm] (clay) at (8.0,0.25)
    {\textbf{Clay} \quad 0\% --- open: the five inputs are unproved};
  \draw[arr] (8.1,3.55) -- (red.north);
  \draw[arr] (11.3,3.55) -- (red.north east);
  \draw[arr] (4.9,3.55) -- (red.north west);
  \draw[arr] (1.7,3.55) -- (red.north west);
  \draw[arr] (14.5,3.55) -- (red.north east);
  \draw[arr] (red.south) -- (clay.north);
\end{tikzpicture}
\caption{What the Lean development does and does not certify. The reduction is machine-checked:
if the five analytic facts hold, then the gap follows, and the proof uses only the standard
axioms. The five facts are hypotheses, not theorems in Lean; three of them are theorems in the
paper (P3, $W_{\mathrm{geo}}\ge0$, no-dissolution), while H1 and (CP) are open. The diagram is
the precise sense in which ``the logic is verified, the inputs are not''.}
\label{fig:proofdag}
\end{figure}
}

% ---------------------------------------------------------------------
% X10 — moments versus determinant
% ---------------------------------------------------------------------
\newcommand{\figfredholm}{%
\begin{figure}[t]
\centering
\begin{tikzpicture}[x=1cm,y=1cm,font=\scriptsize]
  % panel A: the measure
  \draw[arr] (0.9,0.7) -- (6.9,0.7) node[alab, right] {$s$};
  \fill[black!75] (2.4,0.7) rectangle (2.55,2.5);
  \node[above] at (2.47,2.5) {$f_1^2$};
  \fill[gray!25] (3.2,0.7) -- (3.2,1.55) -- (4.2,1.25) -- (5.4,1.0) -- (6.6,0.85) -- (6.6,0.7) -- cycle;
  \node[text=black!60] at (5.0,1.55) {cut};
  \node[text=black!60, align=center, text width=5.2cm] at (3.8,3.65)
    {moments $\int s^k\,d\nu$ see atom and cut alike};
  % panel B: the complex z-plane
  \begin{scope}[shift={(9.4,0)}]
    \draw[arr] (0.2,1.6) -- (5.8,1.6) node[alab, right] {$\Re z$};
    \draw[arr] (3.0,0.35) -- (3.0,2.9) node[alab, above] {$\Im z$};
    \draw[black!45, thick] (3.0,1.6) circle (1.0);
    \node[text=black!55] at (4.55,2.75) {$|z|=1$ (cut)};
    \draw[cN!85, very thick] (4.35,1.6) -- (4.75,1.6);
    \draw[cN!85, very thick] (4.55,1.4) -- (4.55,1.8);
    \node[text=cN!85, below] at (4.55,1.35) {$z_1=e^{am_1}$};
    \node[text=black!60, align=center, text width=5.4cm] at (2.9,3.75)
      {the determinant has a zero at $z_1$ and none elsewhere inside the cut};
  \end{scope}
  \node[text=black!60, align=center, text width=15cm] at (7.7,-0.05)
    {same lower moments, different zero sets: the Fredholm determinant is linear-blind but sees the atom};
\end{tikzpicture}
\caption{Why the Fredholm determinant is the only tool of the six that ``sees'' the atom. Every
moment of the measure --- including all the moments the squeeze uses --- is blind to whether the
bottom of the spectrum is an atom or the edge of a cut. The determinant $\det(1-zT_O)$ is not:
a pole of the resolvent produces a zero at $z_1=e^{am_1}$, a cut produces none. The price is
that the zero has to be followed to $a\to0$.}
\label{fig:fredholm}
\end{figure}
}

% ---------------------------------------------------------------------
% X11 — the verdict card
% ---------------------------------------------------------------------
\newcommand{\figverdictcard}{%
\begin{figure}[t]
\centering
\begin{tikzpicture}[x=1cm,y=1cm,font=\scriptsize]
  \node[fbox, fill=cN!8, draw=cN!60, text width=15.2cm, minimum height=0.9cm] at (7.8,3.9)
    {\textbf{Clay: 0\% --- open.} \ The wall is located and instrumented, not crossed.};
  \foreach \x/\col/\big/\small in {%
    1.9/cT/{100\%}/{G1 assembled\\ modulo $\{L3,G2\}$},
    5.7/cC/{100\%}/{channel pinned\\ (conditional)},
    9.5/cM/{90\% / 0\%}/{G2 region / G2 sign},
    13.3/cM/{6 / 7}/{tools passing\\ acid test / no-gos}}{%
    \node[fbox, fill=\col!7, draw=\col!50, text width=3.3cm, minimum height=2.0cm]
      at (\x,1.9) {\textbf{\big}\\[2pt]\small};}
  \node[text=black!55, align=center, text width=15cm] at (7.8,0.35)
    {the missing object: one single-sided, cutoff-uniform upper bound on a conditioned variance};
\end{tikzpicture}
\caption{The verdict in five numbers. None of them is a percentage of Clay; each is local to its
own object. The two zeros matter more than the three high numbers: G2's sign and the missing
tool are the same absence, and the other numbers only say that everything around it is
assembled.}
\label{fig:verdictcard}
\end{figure}
}

% ---------------------------------------------------------------------
% X12 — the Lean dependency tree
% ---------------------------------------------------------------------
\newcommand{\figlean}{%
\begin{figure}[t]
\centering
\begin{tikzpicture}[x=1cm,y=1cm,font=\scriptsize]
  \node[fbox, fill=cT!10, draw=cT!60, text width=6.8cm] (root) at (8.0,5.5)
    {\code{mass\_gap\_of\_spectral} \quad proved};
  \foreach \x/\nm/\lab/\st in {%
    1.7/{H1\_holds}/{H1}/\LM,
    4.85/{CP\_holds}/{CP}/\LC,
    8.0/{W\_geo\_nonneg\_fact}/{W$_{\mathrm{geo}}\!\ge0$}/\LT,
    11.15/{P3\_strong\_coupling}/{P3}/\LT,
    14.3/{no\_dissolution\_kinematic}/{no-diss.}/\LC}{%
    \node[fbox, fill=cM!8, draw=cM!55, text width=2.8cm, minimum height=1.35cm] at (\x,3.4)
      {\textbf{\lab}\\[1pt]{\tiny\texttt{\nm}}\\[1pt]\textcolor{cM!85}{\st}};
    \draw[arr] (\x,4.1) -- (root.south);}
  \node[fill=black!4, rounded corners=2pt, draw=black!25, text width=11.0cm,
        font=\scriptsize, align=center] (ax) at (8.0,1.35)
    {axioms of \code{clay\_gap\_conditional}:
     \code{[propext, Classical.choice, Quot.sound]} --- no \code{sorryAx}};
\end{tikzpicture}
\caption{The Lean dependency tree. The root lemma is proved; the five leaves are \code{sorry}
hypotheses with the exact analytic statements of the paper. What the machine certifies is the
assembling of the leaves into the gap, not the leaves themselves. This is the formal counterpart
of the label system: a \LT{} inside Lean is a proved lemma, and a \LM{} is a \code{sorry} whose
statement has been fixed.}
\label{fig:lean}
\end{figure}
}
